\documentclass{article}
\usepackage[margin=1in]{geometry}
\usepackage{mathtools, amsfonts, amsthm, xcolor, listings}

\newcommand{\R}{\mathbb{R}}
\newcommand{\N}{\mathbb{N}}
\newcommand{\I}{\mathbb{I}}
\newcommand{\Q}{\mathbb{Q}}
\newcommand{\Z}{\mathbb{Z}}

\newcommand{\set}[1]{\left\{#1\right\}}
\newcommand{\ang}[1]{\left\langle #1 \right\rangle}
\newcommand{\inv}{^{-1}}

\newcommand{\normsub}{\lhd}

\newcommand{\reflection}[5]{
    \begin{itemize}
    \item Points attempted: #1
    \item Self-assessed score: #2/#1.
    \item Time spent: #3 minutes.
    \item Difficulty: #4/10.
    \item Questions/Feedback: #5.
    \end{itemize}
}


\title{HW 11}
\author{Sam Ly}

\begin{document}
\maketitle

\section*{Chapter 22}

\begin{enumerate}
    \item {
        Come up with something more efficient. Do you think the additional complexity is worth it?

        We can use a hashmap to map the variables to their stack index in constant 
        time, however this is not worth it as this cost is only incurred at 
        compile time.

        I would rather have the compiler be simple, and worry about performance 
        when we need to.
    }

    \item {
        What would you do if it was your language? Why?

        Other languages tend to make code like this a compile error. I believe this 
        makes sense, as this is completely ambiguous behavior. One language that 
        I see that allows this is Rust, however Rust semantics allow for this 
        to make sense throw something called `shadowing'.
    }

    \item {
        Pick a keyword for a single-assignment variable form to add to Lox. Justify your choice, then
        implement it. An attempt to assign to a variable declared using your new keyword should
        cause a compile error.

        I would add a keyword `const' that marks a variable as constant. 
        The implementation is too long to fit here, but the overall idea is to 
        add an extra check for the assignment operator, that checks if the 
        left hand side has been marked const.
    }

    \item {
        Extend clox to allow more than 256 local variables to be in scope at a time.

        See code on github.
    }
\end{enumerate}

\section*{Chapter 23}

\begin{enumerate}
    \item {
        To keep things simpler, we’re omitting fallthrough and break statements. Each case
        automatically jumps to the end of the switch statement after its statements are done.

        This can be accomplished through desugaring, since we do not need the 
        fallthrough logic. Switch statements can essentially be converted to 
        if statements.
    }

    \item {
        Make sure to think about scope. What should happen to local variables declared inside the
        body of the loop or in blocks nested inside the loop when a continue is executed?

        When a continue is executed, we need to "reset" the stack to the point 
        at which the loop starts. This way, we don't have access to garbage data,
        and no `stack leaks' occur.
    }

    \item {
        For fun, try to invent a useful novel control flow feature for Lox. It can be a refinement of an
        existing form or something entirely new. In practice, it’s hard to come up with something
        useful enough at this low expressiveness level to outweigh the cost of forcing a user to
        learn an unfamiliar notation and behavior, but it’s a good chance to practice your design
        skills.

        One unique control flow that I have always liked was the `do while' loops 
        from C. Implementing is relatively simple, since this can just be desugared 
        into a normal while loop. 
    }
\end{enumerate}

\end{document}
